\chapter{Project Proceedings}\label{cap:conclusion}

In this thesis proposal, we highlight that despite significant advancements in
image classification, real-world applications like species classification in
camera-trap images still face many challenges. These challenges need to be
addressed to ensure that the data extracted using fully automated AI-based
systems is reliable for use in ecological studies. We presented solutions to
specific challenges to improve classification performance, directly impacting
tasks such as filtering empty camera-trap images, species classification, and
counting individuals. In Section~\ref{sec:contrib}, we present the contributions
of this work, and in Section~\ref{sec:sched}, we outline the next steps of the
research and the schedule.

\section{Contributions}\label{sec:contrib}

The contributions of this work are primarily based on the improvement of
classification models and can be divided according to three tasks in automatic
information extraction from camera-trap images: filtering empty images, species
classification, and counting individuals.

For the task of filtering empty images, we presented a comparative study to
analyze the trade-off between precision and inference latency of classification
and detection models when the inference is done directly on edge devices.
Although some studies have addressed this matter, they typically use small
datasets with a limited number of species, which restricts the scope of their
conclusions. Our work, however, utilizes two large-scale camera-trap datasets to
assess this trade-off more comprehensively. We also applied a series of model
optimizations to reduce inference latency and compared their impact on model
precision. We found that detectors generally outperform classifiers with
comparable inference latency; however, classifiers can achieve higher
performance on this task if a larger amount of labeled images is available,
which is more challenging to obtain for detectors. Another finding is that
models with higher input resolution perform better at filtering empty images,
and their inference latency can be kept competitive by reducing the number of
model filters.

For the widely studied task of species identification, we addressed the problem
of long-tail classification, which is often neglected in works on camera-trap
image species classification, with tail classes typically being removed or
grouped into generic categories. Our main contribution in this area is the
introduction of a new framework called Square-root Sampling Branch (SSB). This
framework employs a multi-branch network combined with sampling strategies to
enhance performance on tail classes while maintaining minimal loss in head
classes' accuracy compared to state-of-the-art long-tail methods. Additionally,
we provide strong classification baselines on four large-scale camera-trap
datasets using the latest training tricks.

For the task of counting animals, we followed the formulation proposed by the
iWildCam 2021 benchmark, where the training set does not contain labels for
counting, necessitating unsupervised learning. Our main contribution was the
development of a strong baseline comprising an ensemble of classifiers
supplemented with counting heuristics based on bounding boxes generated by
object detector models. Our solution won the iWildCam 2021 competition, held at
The Eighth Workshop on Fine-Grained Visual Categorization (FGVC8) during CVPR
2021. Although we explored various methods to obtain more reliable counting
results, such as multi-object tracking and re-identification models based on
bounding boxes, they yielded inferior results compared to our heuristic-based
approach. This highlights the complexity of the problem. In this task, we still
aim to improve the classification subtask by applying the other improvements
proposed in this work.

\section{Next Steps and Schedule}\label{sec:sched}

The next steps of this research will include developing strategies for the 
classification of capture events, applying all classification innovations proposed in
this work to improve the task of counting the number of animals of each species,
and writing the final thesis. \autoref{tab:schedule} shows the schedule of the
activities described as follows:

\begin{itemize}
\item \textbf{A01 - Multi-image experiments}: Develop strategies to improve
species identification based on a sequence of images of the same individual.

\item \textbf{A02 - Multi-image results analysis}: Analyze the results obtained
during the multi-image experiments.

\item \textbf{A03 - Multi-image approach paper}: Write a scientific paper
reporting our proposed multi-image strategies for species classification using a
sequence of images.

\item \textbf{B01 - Improve counting animals task}: Apply all classification
innovations proposed in this work to improve the task of counting the number of
animals of each species.

\item \textbf{B02 - Counting animals results analysis}: Analyze the results on
the task of classifying species and counting individual animals present in a
sequence of images.

\item \textbf{C01 - Thesis writing}: Write the final thesis.

\item \textbf{C02 - Thesis defense}
\end{itemize}



\begin{table}[h]
\centering
\caption{Research schedule by month}
\label{tab:schedule}
\begin{tabular}{|c|lllllll|}
\hline
\multicolumn{1}{|c|}{} &
  \multicolumn{7}{c|}{2025} \\ \cline{2-8}
\multicolumn{1}{|c|}{\multirow{-2}{*}{Activities}} &
  \multicolumn{1}{c|}{Jun} &
  \multicolumn{1}{c|}{Jul} &
  \multicolumn{1}{c|}{Ago} &
  \multicolumn{1}{c|}{Sep} &
  \multicolumn{1}{c|}{Oct} &
  \multicolumn{1}{c|}{Nov} &
  \multicolumn{1}{c|}{Dec} \\ \hline
A01 &
  \multicolumn{1}{l|}{\cellcolor[HTML]{9B9B9B}{\color[HTML]{330001} }} &
  \multicolumn{1}{l|}{\cellcolor[HTML]{9B9B9B}{\color[HTML]{330001} }} &
  \multicolumn{1}{l|}{} &
  \multicolumn{1}{l|}{} &
  \multicolumn{1}{l|}{} &
  \multicolumn{1}{l|}{} &
   \\ \hline
A02 &
  \multicolumn{1}{l|}{} &
  \multicolumn{1}{l|}{\cellcolor[HTML]{9B9B9B}} &
  \multicolumn{1}{l|}{} &
  \multicolumn{1}{l|}{} &
  \multicolumn{1}{l|}{} &
  \multicolumn{1}{l|}{} &
   \\ \hline
A03 &
  \multicolumn{1}{l|}{} &
  \multicolumn{1}{l|}{\cellcolor[HTML]{9B9B9B}} &
  \multicolumn{1}{l|}{\cellcolor[HTML]{9B9B9B}} &
  \multicolumn{1}{l|}{\cellcolor[HTML]{9B9B9B}} &
  \multicolumn{1}{l|}{} &
  \multicolumn{1}{l|}{} &
   \\ \hline
B01 &
  \multicolumn{1}{l|}{} &
  \multicolumn{1}{l|}{} &
  \multicolumn{1}{l|}{} &
  \multicolumn{1}{l|}{\cellcolor[HTML]{9B9B9B}} &
  \multicolumn{1}{l|}{} &
  \multicolumn{1}{l|}{} &
   \\ \hline
B02 &
  \multicolumn{1}{l|}{} &
  \multicolumn{1}{l|}{} &
  \multicolumn{1}{l|}{} &
  \multicolumn{1}{l|}{} &
  \multicolumn{1}{l|}{\cellcolor[HTML]{9B9B9B}} &
  \multicolumn{1}{l|}{} &
   \\ \hline
C01 &
  \multicolumn{1}{l|}{} &
  \multicolumn{1}{l|}{} &
  \multicolumn{1}{l|}{} &
  \multicolumn{1}{l|}{\cellcolor[HTML]{9B9B9B}} &
  \multicolumn{1}{l|}{\cellcolor[HTML]{9B9B9B}} &
  \multicolumn{1}{l|}{\cellcolor[HTML]{9B9B9B}} & 
  \multicolumn{1}{l|}{} \\ \hline
C02 &
  \multicolumn{1}{l|}{} &
  \multicolumn{1}{l|}{} &
  \multicolumn{1}{l|}{} &
  \multicolumn{1}{l|}{} &
  \multicolumn{1}{l|}{} &
  \multicolumn{1}{l|}{} &
  \multicolumn{1}{l|}{\cellcolor[HTML]{9B9B9B}} \\ \hline
\end{tabular}
\end{table}
